%%% FIELD proof-local-contrast
For \(r\in\{0,1\}\), let \(\Phi^{(r)}(y)\) denote the Eq.\ 23 expression obtained by putting \(q_C^{(r)}\) in the \(C\)-factor and retaining the common outer conditional \(p^S\) and the common factors \(q_j\), \(D_j\ne C\).  By \(\mathrm{def:eq23\text{-}functional}\), all state spaces and all displayed sums are finite.  Hence ordinary distributivity permits termwise subtraction, and for every \(y\) one has
\[
\begin{aligned}
 \Phi^{(0)}(y)-\Phi^{(1)}(y)
 &=\sum_{w,d_{D\setminus Y_N}}p^S(y_S,w\mid x_S)
 \left\{q_C^{(0)}\prod_{D_j\ne C}q_j-q_C^{(1)}\prod_{D_j\ne C}q_j\right\}\\
 &=\sum_{w,d_{D\setminus Y_N}}p^S(y_S,w\mid x_S)
 \bigl(q_C^{(0)}-q_C^{(1)}\bigr)\prod_{D_j\ne C}q_j\\
 &=\Delta(y).
\end{aligned}
\]
The first equality is substitution in the defining product, the second is distributivity, and the third is the definition of \(\Delta\) in the statement.  Two probability measures on the finite state space \(\prod_{v\in Y}\mathcal X_v\) are unequal if and only if their masses differ at some \(y\).  Therefore the two constructed targets differ if and only if \(\Delta(y)\ne0\) for at least one \(y\).

It remains to verify the claimed obstruction.  Let a local factor have two binary coordinates \((A,R)\), and define
\[
\begin{array}{c|cccc}
 (a,r)&(0,0)&(1,0)&(0,1)&(1,1)\\ \hline
 q_C^{(0)}(a,r)&3/8&1/8&1/8&3/8\\
 q_C^{(1)}(a,r)&1/4&1/4&1/4&1/4.
\end{array}
\]
Both rows are strictly positive probability masses and they are unequal.  Nevertheless, for each \(r\in\{0,1\}\),
\[
 \sum_{a\in\{0,1\}}\bigl(q_C^{(0)}(a,r)-q_C^{(1)}(a,r)\bigr)=0.
\]
Place \(A\) among the coordinates summed in \(D\setminus Y_N\), and choose the common outer conditional and all common factors to be strictly positive and constant in \(a\).  Such a choice is allowed in the algebraic assertion, and the empty product is interpreted as one if there is no other district.  The sum defining \(\Delta(y)\) then contains, for each fixed value of its remaining indices, a positive common multiplier times the last displayed zero.  Thus \(\Delta(y)=0\) for every \(y\), although every common factor is positive and \(q_C^{(0)}\ne q_C^{(1)}\).  This proves both the equivalence and the noninjectivity assertion.

%%% FIELD current-district-extension
Fix \(D_i\in\mathcal D_S(D)\) and its initial enclosure \(T_i^{(0)}\in\operatorname{SComp}_{G^S}(N_S)\) in the exact induced graph \(G^S\). For a call chain with fixed target \(C=D_i\) that ends at a terminal failure \(C\subsetneq T\subseteq T_i^{(0)}\) satisfying \(\operatorname{SHedge}_{G^S}(C,T)\), let \((F_0,F_1,R,\overline M_0,\overline M_1,x_{\mathrm{bin}})\in\mathsf{ProjHedgeInput}(G^S;X,Y;C,T)\). Any local exact rational pair supplied by the effective \(\mathsf H_{\mathrm{lift}}\) for this record extends through the selected district factorization to \(\mathsf{SHedgePair}(G^S,C,T,x_{\mathrm{bin}};\sigma_{\mathrm{bin}})=(E_0,E_1)\in\operatorname{CertSCM}(G^S,\sigma_{\mathrm{bin}})^2\). The extension appends identical rational strictly positive kernels outside the failed component; \(\mathrm{lem:local\text{-}contrast\text{-}survives\text{-}eq23\) turns the preserved \(D_i=C\) kernel contrast into a nonzero contrast in the original \(\Theta_{x_{\mathrm{bin}},\sigma_{\mathrm{bin}}}\); and \(\operatorname{Exactify}_{G^S}\) completes every parent interface and shared-latent incidence without changing any selected atom or interventional target. It yields \(\mathsf{Check}_{\mathrm{NONID}}(G^S,X,Y,x_{\mathrm{bin}},E_0,E_1)=\mathrm{accept}\); consequently both realized models induce exactly \(G^S\), have one positive selected law, and have unequal targets.

%%% FIELD proposed-district-extension
Fix \(D_i\in\mathcal D_S(D)\) and its initial enclosure \(T_i^{(0)}\in\operatorname{SComp}_{G^S}(N_S)\).  Suppose a terminal call with fixed target \(C=D_i\) fails at \(C\subsetneq T\subseteq T_i^{(0)}\), and fix \((F_0,F_1,R,\overline M_0,\overline M_1,x_{\mathrm{bin}})\in\mathsf{ProjHedgeInput}(G^S;X,Y;C,T)\).  Consider any two finite rational subgraph-compatible SCM encodings that extend the projected pair and have the same strictly positive selected observational kernel.  Suppose their postintervention target masses have the common-factor form in Eq.\ 23: the two products use the same outer selected conditional \(p^S(y_S,w\mid x_S)>0\) and the same strictly positive factors \(q_j\) for \(D_j\ne C\), and use factors \(q_C^{(0)}\) and \(q_C^{(1)}\) in the \(C\)-position.  For every outcome value \(y\), define
\[
 \Delta(y)=\sum_{w,d_{D\setminus Y_N}}p^S(y_S,w\mid x_S)
 \bigl(q_C^{(0)}-q_C^{(1)}\bigr)\prod_{D_j\ne C}q_j.
\]
After applying \(\operatorname{Exactify}_{G^S}\), the two encodings \(E_0,E_1\) belong to \(\operatorname{CertSCM}(G^S,\sigma_{\mathrm{bin}})\), induce exactly \(G^S\), and retain the same strictly positive selected law.  For this fixed common-factor extension,
\[
 \mathsf{Check}_{\mathrm{NONID}}(G^S,X,Y,x_{\mathrm{bin}},E_0,E_1)=\mathrm{accept}
 \quad\Longleftrightarrow\quad
 \Delta(y)\ne0\text{ for some }y.
\]
When these equivalent conditions hold, the pair is the certified \(\mathsf{SHedgePair}(G^S,C,T,x_{\mathrm{bin}};\sigma_{\mathrm{bin}})\).  Local kernel inequality and strict positivity alone do not imply either equivalent condition.

%%% FIELD proof-district-extension
Work first before exactification.  By hypothesis the two candidate SCM extensions have identical strictly positive selected observational kernels.  In their common-factor representation, every appended outer conditional and every observational district factor is identical in the two encodings.  The selected district factorization in \(\mathrm{def:selected\text{-}kernel}\) therefore gives, cell by cell,
\[
 p^S\,q_{C,\mathrm{obs}}\prod_{D_j\ne C}q_j
 =p^S\,q_{C,\mathrm{obs}}\prod_{D_j\ne C}q_j.
\]
Every factor is positive, so every selected cell in the common product is positive.  The assumed finite rational subgraph-compatible SCM encodings are therefore elements of \(\operatorname{RawCertSCM}(G^S,\sigma_{\mathrm{bin}})\) with the same strictly positive selected law.

Under \(\operatorname{do}(X=x_{\mathrm{bin}})\), the only factor allowed to differ between the two displayed Eq.\ 23 products is \(q_C^{(r)}\).  By \(\mathrm{lem:local\text{-}contrast\text{-}survives\text{-}eq23}\), their point-mass difference at \(y\) is exactly \(\Delta(y)\).

Apply \(\operatorname{Exactify}_{G^S}\) to both pre-exactification encodings.  By \(\mathrm{lem:exactification\text{-}preserves\text{-}laws}\), the results \(E_0,E_1\) belong to \(\operatorname{CertSCM}(G^S,\sigma_{\mathrm{bin}})\), induce exactly the required parent interfaces and bidirected incidences, and preserve the common selected law and both interventional targets.  By \(\mathrm{lem:certscm\text{-}ambient\text{-}inclusion}\) and \(\mathrm{lem:target\text{-}support}\), the two postintervention conditioning probabilities are positive.  Thus all structural, rationality, positivity, and selected-law checks in \(\mathrm{def:certificate\text{-}checker}\) pass.  Its remaining acceptance test is precisely inequality of the two finite target distributions.

Keep the common extension fixed.  The subtraction in \(\mathrm{lem:local\text{-}contrast\text{-}survives\text{-}eq23}\) shows that the target point masses differ exactly by \(\Delta(y)\).  Since the outcome space is finite, the target distributions are unequal exactly when \(\Delta(y)\ne0\) for some outcome value.  By the preceding description of the checker, this is also equivalent to acceptance.  When acceptance holds, \(\mathrm{lem:certificate\text{-}checker\text{-}soundness}\) supplies the advertised exact-graph countermodel consequences, so the accepted pair has the defining evidence required of \(\mathsf{SHedgePair}\).

Finally, the positive binary tables in \(\mathrm{lem:local\text{-}contrast\text{-}survives\text{-}eq23}\) have \(q_C^{(0)}\ne q_C^{(1)}\) but \(\Delta(y)=0\) after a marginalized coordinate is summed.  They prove that local inequality and positivity alone do not force acceptance.  The statement is therefore a characterization of a candidate extension, not an existence claim for the unresolved compiler.

%%% FIELD graphwise-open-what
The failure-to-nonidentification implication is open: one must construct, for every returned \(S\)-Hedge with shared or branching paths through \(A_S\cup\{S\}\), a finite rational exact-graph pair with a common positive selected law and with a nonzero full-target Eq.\ 23 contraction.

%%% FIELD graphwise-open-obstruction
The projected hedge supplies only \(q_C^{(0)}\ne q_C^{(1)}\).  The proved contraction lemma exhibits a nontrivial nullspace, so this inequality does not imply separation of \(\Theta_{x_{\mathrm{bin}},\sigma_{\mathrm{bin}}}\).  The frozen core also contains no proved joint inverse lift for arbitrary overlapping ancestor paths and no proved positive rational compiler for the initial selected d-separation failure.  Construction definitions that name \(\mathsf H_{\mathrm{lift}}\), \(\mathsf{SHedgePair}\), and \(\mathsf{AncestorPair}\) do not prove those properties.

%%% FIELD graphwise-open-attempted
The terminal trace was reduced to ordinary hedge geometry through the suppression projection, and the exact two-virtual-edge graph was checked as a positive binary base case.  A common-factor district extension was then expanded equation by equation.  This calculation isolates the missing condition as \(\Delta(y)\ne0\); direct positive binary kernels show that marginalization can erase the local contrast, so positivity and ordinary hedge separation alone cannot finish the lift.

%%% FIELD graphwise-open-partial
The success direction is complete.  By \(\mathrm{lem:published\text{-}latent\text{-}sid\text{-}soundness}\), whenever Algorithm 1 succeeds, for every finite signature, every compatible positive model, and every intervention value, \(\Phi_{23,x}^{\sigma}(P_{M,\sigma}^S)=\Theta_{x,\sigma}(M)\).  Hence equal selected laws give equal targets.  On a recursive failure, \(\mathrm{lem:shedge\text{-}projects\text{-}to\text{-}hedge}\) gives the projected nested forests.  Conditional on a rational inverse lift and on a common-factor extension satisfying the computable condition \(\Delta(y)\ne0\) for some \(y\), the corrected district-extension lemma and checker soundness give two members of \(\mathfrak M_+(G^S,\sigma_{\mathrm{bin}})\) with one positive selected law and unequal targets.  The missing result is the uniform construction of such a contrast-transmitting lift, together with the separate initial-failure compiler.

%%% FIELD current-executable
Assume the joint fiber-coupling handle is proved to be an effective constructor returning rational \(\operatorname{CertSCM}\) pairs with accepted checker evidence for every actual terminal Algorithm 1 failure. Then, for every finite augmented ADMG \(G^S\) with childless \(S\), interpreted as the exact graph induced by declared parent sets and shared exogenous inputs, every finite signature \(\sigma\), and every \(x\in\prod_{v\in X}\mathcal X_v\), \(\mathsf{sID}^{\pm}\) is total. It returns either \(\operatorname{ID}(\Phi_{23,x}^{\sigma})\) or \(\operatorname{NONID}(x_{\mathrm{bin}},E_0,E_1,\mathsf{Check}_{\mathrm{NONID}}=\mathrm{accept};\sigma_{\mathrm{bin}})\). Its initial call family is \(\mathcal I_0=\{(D_i,T_i^{(0)}):D_i\in\mathcal D_S(D)\}\). Along the chain rooted at \((D_i,T_i^{(0)})\), the target remains \(C=D_i\); a recursive step forms \(A=\operatorname{An}_{G^S[T]}(D_i)\) and replaces \(T\) by the S-component \(T'\subseteq A\subsetneq T\) containing \(D_i\). Thus \(\rho\) decreases strictly. The chain rooted at \(T_i^{(0)}\) contains at most \(|T_i^{(0)}|\) single-component invocations, the recursion depth is at most \(\max_i|T_i^{(0)}|\le|N_S|\), and the total number of single-component invocations is at most \(B_0=\sum_{(D_i,T_i^{(0)})\in\mathcal I_0}|T_i^{(0)}|\). For each returned NONID pair, the checker validates exact induced-graph equality and terminates within \(c\{|V|^2+L(E_0)+L(E_1)+2^{|V|}(R(E_0)+R(E_1))(|V|+|Y|+1)\}\) rational arithmetic and table-lookup operations; acceptance certifies equal positive selected laws and unequal targets. An existential witness in an arbitrary SCM representation does not establish this totality claim.

%%% FIELD proposed-executable
Assume that the initial selected d-separation check of published Algorithm 1 succeeds.  Assume also that, for every terminal single-component failure reached by this run, an effective constructor returns \(x_{\mathrm{bin}}\in\{0,1\}^{X}\) and \(E_0,E_1\in\operatorname{CertSCM}(G^S,\sigma_{\mathrm{bin}})\) with \(\mathsf{Check}_{\mathrm{NONID}}(G^S,X,Y,x_{\mathrm{bin}},E_0,E_1)=\mathrm{accept}\).  Then, for every finite augmented ADMG \(G^S\) with childless \(S\), every finite signature \(\sigma\), every \(M\in\mathfrak M_+(G^S,\sigma)\), and every \(x\in\prod_{v\in X}\mathcal X_v\), the corresponding restricted run of \(\mathsf{sID}^{\pm}\) on \(P_{M,\sigma}^S\) is total.  It returns either \(\operatorname{ID}(\Phi_{23,x}^{\sigma})\), in which case \(\Phi_{23,x}^{\sigma}(P_{M,\sigma}^S)=\Theta_{x,\sigma}(M)\), or the stated \(\operatorname{NONID}\) object.  Its initial call family is \(\mathcal I_0=\{(D_i,T_i^{(0)}):D_i\in\mathcal D_S(D)\}\).  Along each chain the target remains \(C=D_i\), and every recursive step replaces an enclosure \(T\) by the S-component \(T'\) containing \(D_i\) inside \(\operatorname{An}_{G^S[T]}(D_i)\subsetneq T\).  Consequently \(\rho\) decreases strictly, that chain has at most \(|T_i^{(0)}|\) invocations, the recursion depth is at most \(\max_i|T_i^{(0)}|\le|N_S|\), and the total number of single-component invocations is at most \(B_0=\sum_{(D_i,T_i^{(0)})\in\mathcal I_0}|T_i^{(0)}|\).  For every returned NONID pair, the checker terminates within \(c\{|V|^2+L(E_0)+L(E_1)+2^{|V|}(R(E_0)+R(E_1))(|V|+|Y|+1)\}\) exact rational arithmetic and table-lookup operations, and acceptance certifies exact induced-graph equality, equal positive selected laws, and unequal targets.  This result makes no claim about a run whose initial selected d-separation check fails.

%%% FIELD proof-executable
Because the initial selected d-separation check succeeds, the restricted run reaches the finite ancestral set \(D\), its finite S-component partition \(\mathcal D_S(D)\), and the initial family \(\mathcal I_0\) described in \(\mathrm{def:initial\text{-}call\text{-}budget}\).  Fix a root \((D_i,T_i^{(0)})\).  By \(\mathrm{def:success\text{-}or\text{-}certificate\text{-}algorithm}\), the first argument of every single-component call on this chain remains \(C=D_i\).  A recursive call forms
\[
 A=\operatorname{An}_{G^S[T]}(D_i)
\]
and is made only in the strict-shrinkage case; its new enclosure is the S-component \(T'\) of \(A\) containing \(D_i\).  Therefore
\[
 \varnothing\ne D_i\subseteq T'\subseteq A\subsetneq T,
 \qquad |T'|\le |T|-1.
\]
This proves directly that the multiset rank \(\rho\) of \(\mathrm{def:termination\text{-}rank}\) decreases at every recursive call.

If \(|T_i^{(0)}|=m\), its enclosure size is a positive integer initially equal to \(m\) and decreases by at least one at each recursion.  There can therefore be at most \(m-1\) recursive transitions and at most \(m\) invocations including the initial call.  Since there are finitely many roots, every chain terminates.  Taking the maximum over roots and then summing their individual bounds gives
\[
 \operatorname{depth}\le\max_i|T_i^{(0)}|\le |N_S|
\]
and
\[
 \#\{\text{single-component invocations}\}
 \le \sum_{(D_i,T_i^{(0)})\in\mathcal I_0}|T_i^{(0)}|
 =B_0,
\]
where the last equality is \(\mathrm{def:initial\text{-}call\text{-}budget}\).  This elementary size argument also proves termination without invoking an undeclared well-foundedness theorem.

At termination there are two cases.  If all component calls succeed, \(\mathsf{sID}^{\pm}\) returns the Eq.\ 23 functional by \(\mathrm{def:success\text{-}or\text{-}certificate\text{-}algorithm}\).  The cited soundness result \(\mathrm{lem:published\text{-}latent\text{-}sid\text{-}soundness}\) yields
\[
 \Phi_{23,x}^{\sigma}(P_{M,\sigma}^S)=\Theta_{x,\sigma}(M).
\]
The displayed selected-law conditionals have positive denominators because every selected cell is positive by \(M\in\mathfrak M_+(G^S,\sigma)\), and \(\mathrm{lem:target\text{-}support}\) gives \(P_M(S=1\mid\operatorname{do}(X=x))>0\).  Thus both sides are defined.

Otherwise a terminal single-component failure occurs.  The second premise supplies the finite rational intervention value and certificates, so the algorithm returns the stated NONID object.  By \(\mathrm{lem:certificate\text{-}checker\text{-}soundness}\), acceptance implies that both realizations lie in \(\mathfrak M_+(G^S,\sigma_{\mathrm{bin}})\), induce exactly \(G^S\), have the same strictly positive selected law, and have unequal intervention-before-conditioning targets.  The same lemma gives the operation bound
\[
 c\left\{|V|^2+L(E_0)+L(E_1)
 +2^{|V|}\bigl(R(E_0)+R(E_1)\bigr)(|V|+|Y|+1)\right\}.
\]
These cases are exhaustive for a run whose initial check has succeeded, proving totality and all asserted guarantees.  No property of \(\mathsf{AncestorPair}\) is used, which is why the conclusion is explicitly restricted away from the unresolved initial-failure branch.
